\chapter{Approche algorithmique séquentielle}
\label{chap:approche}

\epigraph{\textit{``The art of programming is the art of organizing complexity.''}}{--- Edsger W. Dijkstra}

Ce chapitre détaille l'algorithme WFC tel que nous l'avons implémenté en série. Il sert de référence pour les chapitres parallèles : les versions OpenMP et Kokkos partagent exactement la même logique, seuls les schémas d'orchestration diffèrent.

\section{Les six étapes de WFC}

L'algorithme WFC se résume en six étapes, alignées sur la formulation de Gumin \cite{gumin2016wfc} et la présentation pédagogique de Heaton \cite{heaton2018wfc}.

\begin{algorithm}[H]
\caption{Wave Function Collapse séquentiel}
\begin{algorithmic}[1]
\Require Échantillon $S$, taille de tuile $N$, dimensions sortie $(r, c)$, seed
\Ensure Grille $G$ ou échec
\State $T, \text{freq} \gets \textsc{ExtraireTuiles}(S, N)$
\State $R \gets \textsc{ConstruireRègles}(T, N)$
\State $W \gets \textsc{InitWave}(r, c, |T|)$ \Comment{toutes tuiles candidates partout}
\State $\text{rng} \gets \textsc{Mt19937}(\text{seed})$
\Loop
    \State $\text{cell} \gets \arg\min_c \textsc{Entropy}(W[c]) + \textsc{Jitter}(c, \text{seed})$
    \If{$\text{cell} = -1$}
        \State \Return \textsc{BuildOutput}($W$, $T$) \Comment{succès}
    \EndIf
    \State $t \gets \textsc{WeightedPick}(W[\text{cell}], \text{freq}, \text{rng})$
    \State $W[\text{cell}] \gets \{t\}$ \Comment{collapse}
    \If{$\neg \textsc{Propagate}(W, R, \text{cell})$}
        \State \Return $\bot$ \Comment{contradiction}
    \EndIf
\EndLoop
\end{algorithmic}
\end{algorithm}

\subsection{Étape 1 : Extraction des tuiles}

\begin{definition}[Extraction toroïdale]
Pour chaque position $(r_0, c_0) \in [0, r_S) \times [0, c_S)$ de l'échantillon, on lit $N \times N$ valeurs avec wrap-around :
\[
\text{tile}[i][j] = S[(r_0 + i) \bmod r_S][(c_0 + j) \bmod c_S]
\]
\end{definition}

L'algorithme accumule les tuiles dans une \texttt{unordered\_map<Tile, int>} pour dédupliquer en $O(1)$ amorti par hash FNV-1a. Pour chaque tuile rencontrée, sa fréquence est incrémentée. Le résultat : un vecteur $T$ des tuiles uniques et un vecteur \texttt{freq} parallèle.

\paragraph{Choix toroïdal vs non-toroïdal.} Le sujet ne précise pas la convention. Nous choisissons \textbf{toroïdal} pour deux raisons :
\begin{itemize}
    \item Chaque pixel apparaît exactement dans $N \times N$ tuiles, distribution uniforme.
    \item Le nombre total de tuiles extraites vaut $r_S \times c_S$ (et non $(r_S - N + 1) \times (c_S - N + 1)$), ce qui évite les biais de bord.
\end{itemize}

C'est aussi le choix de l'implémentation référence de Gumin. Pour l'exemple du sujet ($5 \times 5$ binaire avec $N = 2$), l'extraction toroïdale identifie 11 tuiles uniques (vs 8 en non-toroïdal). Tests \texttt{test\_tileset} valident ces 11 tuiles bit-à-bit.

\subsection{Étape 2 : Pré-calcul des règles d'adjacence}

Pour chaque paire de tuiles $(t_1, t_2)$ et chaque offset $(dx, dy) \in [-(N-1), N-1]^2$, on calcule un booléen : « $t_2$ est-elle compatible avec $t_1$ au décalage $(dx, dy)$ ? ». Le résultat est stocké comme \textbf{bitset} indexé par identifiant de tuile.

\begin{figure}[H]
\centering
\schemaImage{Compatibilite overlap entre tuiles 2x2}
\caption{Convention d'overlap : $t_2$ placée à $(dx, dy)$ de $t_1$, les valeurs doivent coïncider sur la zone de recouvrement.}
\label{fig:overlap_compatibility}
\end{figure}

Le stockage est plat : un \texttt{vector<Bitset>} de taille $L \cdot (2N-1)^2$, indexé par
\[
\text{rule\_idx}(t, dx, dy) = t \cdot (2N-1)^2 + (dx + N - 1) \cdot (2N-1) + (dy + N - 1)
\]

\paragraph{Coût.} $L^2 \times (2N-1)^2$ paires à tester, chaque test en $O(N^2)$. Pour le sample du sujet ($L = 11$, $N = 2$) : $\sim 5000$ comparaisons, exécutées en quelques µs. Pour $L = 73$, $N = 3$ : $\sim 500\,000$ comparaisons, encore quelques ms. Cette construction est parallélisée en OpenMP via \texttt{\#pragma omp parallel for} sur la boucle externe (gain $\sim 2\times$ sur 8 cores).

\subsection{Étape 3 : Initialisation de la Wave}

Chaque cellule $(r, c)$ de la grille de sortie reçoit le bitset \textbf{plein} (toutes les $L$ tuiles candidates). Le constructeur de \texttt{Wave} fait un \textit{first-touch parallèle} (\texttt{\#pragma omp parallel for} sur l'init) pour mapper les pages physiques sur le nœud NUMA des threads consommateurs.

\subsection{Étape 4 : Sélection min-entropie}

\begin{definition}[Entropie de Shannon pondérée]
Pour une cellule $c$ dont l'ensemble de candidats restants est $W[c]$, l'entropie est :
\[
H(c) = \log\left(\sum_{t \in W[c]} f_t\right) - \frac{\sum_{t \in W[c]} f_t \log f_t}{\sum_{t \in W[c]} f_t}
\]
où $f_t$ est la fréquence de la tuile $t$ dans $S$.
\end{definition}

L'algorithme choisit la cellule de moindre entropie (cellule la plus contrainte). Pour départager les ex-aequo (très fréquents en début de résolution quand toutes les cellules ont la même entropie), on ajoute un \textbf{bruit déterministe} :
\[
H'(c) = H(c) + \text{jitter}(c, \text{seed})
\]
où \texttt{jitter} est un hash SplitMix64 de $(c, \text{seed})$ produisant une valeur dans $[0, 10^{-6})$. Ce bruit est \textbf{stateless} : il ne consomme pas l'état du RNG, ce qui garantit le déterminisme à travers les backends parallèles (voir Chapitre~\ref{chap:openmp}).

\subsection{Étape 5 : Tirage pondéré et collapse}

Une fois la cellule choisie, on tire une tuile selon les fréquences :
\begin{enumerate}
    \item Calculer la somme $S = \sum_{t \in W[\text{cell}]} f_t$
    \item Tirer $u \sim \mathcal{U}(0, S)$
    \item Parcourir les candidats par ordre croissant de $t$, accumuler la CDF, retourner le premier $t$ tel que $\text{CDF}(t) \geq u$.
\end{enumerate}

\texttt{wave[cell].set\_only(t)} réduit le bitset à un seul bit (force la collapse).

\subsection{Étape 6 : Propagation BFS}

\begin{algorithm}[H]
\caption{Propagation BFS niveau-synchrone}
\begin{algorithmic}[1]
\Require Wave $W$, règles $R$, cellule de départ $c_0$
\Ensure $\bot$ si contradiction, $\top$ sinon
\State $Q \gets \{c_0\}$
\While{$Q \neq \emptyset$}
    \State $c \gets Q.\textsc{Pop}()$
    \For{chaque offset $(dx, dy) \in [-(N-1), N-1]^2 \setminus \{(0,0)\}$}
        \State $nc \gets \textsc{Voisin}(c, dx, dy)$
        \State $A \gets \bigcup_{t \in W[c]} R(t, dx, dy)$ \Comment{union des compatibilités}
        \State $\text{changed} \gets W[nc] \gets W[nc] \cap A$ \Comment{intersection bitwise}
        \If{$\text{changed}$}
            \If{$W[nc] = \emptyset$}
                \State \Return $\bot$ \Comment{contradiction}
            \EndIf
            \State $Q.\textsc{Push}(nc)$
        \EndIf
    \EndFor
\EndWhile
\State \Return $\top$
\end{algorithmic}
\end{algorithm}

\paragraph{Variante \texttt{for\_each\_set} optimisée.} L'union $A$ se calcule en énumérant les bits actifs de $W[c]$ via \texttt{\_\_builtin\_ctzll} (count trailing zeros). Pour chaque bit actif, on \texttt{or\_with} la règle correspondante. Beaucoup plus rapide que tester chaque bit en $O(L)$.

% =============================================================================
\section{Stratégies de récupération sur contradiction}
\label{sec:approche:retry}
% =============================================================================

\subsection{Restart-on-contradiction (défaut)}

Quand \texttt{propagate} retourne $\bot$, l'algorithme retourne \texttt{false} pour le \texttt{run\_attempt} courant. La boucle \texttt{solve()} retente avec un nouveau seed dérivé :
\[
\text{seed}_k = \text{seed}_{\text{base}} + k \cdot \phi
\]
où $\phi = \texttt{0x9E3779B97F4A7C15}$ est l'inverse du nombre d'or 64 bits (Knuth). Jusqu'à \texttt{max\_attempts} essais.

\paragraph{Pourquoi ce choix ?} Trois raisons :
\begin{itemize}
    \item Backtracking propre = sauvegarder/restaurer l'état complet de la wave, potentiellement profond, avec impact mémoire imprédictible.
    \item En pratique, sur les samples fournis, 1-2 attempts suffisent.
    \item Pédagogiquement plus simple et lisible.
\end{itemize}

\subsection{Backtracking (option \texttt{--backtrack})}

L'extension \texttt{--backtrack} (Chapitre~\ref{chap:extensions}) implémente l'alternative : sur contradiction, dépiler le dernier point de décision et essayer le candidat suivant. Le snapshot delta-encodé minimise la mémoire à $\sim 50$ cellules par frame plutôt que la wave complète.

% =============================================================================
\section{Reconstruction de la sortie}
\label{sec:approche:output}
% =============================================================================

Une fois la wave complètement collapsée (chaque cellule à 1 candidat), on reconstruit la grille $G$ :

\begin{definition}[Reconstruction]
Pour chaque cellule $(r, c)$ de la wave, on lit l'unique tuile $t = W[r][c]$ encore active. La valeur de pixel $G[r][c]$ est la valeur du \textbf{coin haut-gauche} de $t$ :
\[
G[r][c] = T[t][0][0]
\]
\end{definition}

Cette convention est nécessaire car dans le modèle overlapping, plusieurs tuiles couvrent chaque pixel. On choisit conventionnellement celui dont le coin haut-gauche est sur $(r, c)$. Vérifié par \texttt{test\_solver::output\_uses\_only\_input\_tiles}.

% =============================================================================
\section{Complexité et profilage}
\label{sec:approche:complexite}
% =============================================================================

\begin{table}[H]
\centering
\begin{tabular}{lll}
\toprule
\textbf{Étape} & \textbf{Complexité} & \textbf{Pour 128×128, $L = 11$} \\
\midrule
Extraction tuiles & $O(r_S \cdot c_S \cdot N^2)$ & $\sim 100$ µs (une fois) \\
Règles d'adjacence & $O(L^2 \cdot (2N-1)^2 \cdot N^2)$ & $\sim 50$ µs (une fois) \\
\textbf{Min-entropie} & $O(r \cdot c \cdot L)$ \textbf{par collapse} & $\sim 30$ µs \textbf{par collapse} \\
Propagation BFS & $O(r \cdot c \cdot (2N-1)^2 \cdot L)$ amorti par collapse & $\sim 150$ µs par collapse \\
Tirage RNG & $O(L)$ par collapse & $\sim 1$ µs par collapse \\
\midrule
\textbf{Total} (8000 collapses) & dominé par min-entropie + BFS & $\sim 3-4$ s en série \\
\bottomrule
\end{tabular}
\caption{Décomposition de complexité du WFC séquentiel sur grille 128×128 binaire ($L = 11$).}
\label{tab:wfc_complexity}
\end{table}

\begin{important}{Résultat contre-intuitif}
La sélection min-entropie (re-scan de la wave entière à chaque collapse) domine $\sim 85\%$ du temps total. La propagation BFS, malgré sa complexité visible plus grande, ne représente que $\sim 12\%$.

C'est ce qui justifie l'effort de parallélisation principal sur \texttt{parallel\_min\_entropy}, et non sur \texttt{propagate\_tasks} (Chapitre~\ref{chap:openmp}).
\end{important}

\section{Garanties d'invariants}

Notre suite de tests vérifie systématiquement plusieurs invariants algorithmiques (Tableau~\ref{tab:invariants}).

\begin{table}[H]
\centering
\begin{tabularx}{\textwidth}{lX}
\toprule
\textbf{Invariant} & \textbf{Vérifié par} \\
\midrule
Sortie utilise uniquement des tuiles présentes dans l'échantillon & \texttt{test\_solver::output\_uses\_only\_input\_tiles} \\
\midrule
Symétrie des règles : $t_2 \in \text{allow}(t_1, d) \iff t_1 \in \text{allow}(t_2, -d)$ & \texttt{test\_overlap} \\
\midrule
Identité au décalage $(0, 0)$ : $\text{allow}(t, 0, 0) = \{t\}$ & \texttt{test\_overlap} \\
\midrule
Total des fréquences = $r_S \cdot c_S$ (toroïdal) & \texttt{test\_tileset} \\
\midrule
Déterminisme : même seed → même output (bit-à-bit) & \texttt{test\_solver}, \texttt{test\_solver\_omp}, \texttt{test\_solver\_kokkos} \\
\midrule
Bit-identité serial = OMP = Kokkos pour un même seed et $\{1, 2, 4, 8\}$ threads & \texttt{test\_solver\_omp}, \texttt{test\_solver\_kokkos} \\
\midrule
Reseed produit une sortie différente (haute probabilité) & \texttt{test\_solver\_omp}, \texttt{test\_solver\_kokkos} \\
\bottomrule
\end{tabularx}
\caption{Invariants algorithmiques vérifiés par la suite de tests.}
\label{tab:invariants}
\end{table}
